\documentclass[12pt,english]{article}
\usepackage[T1]{fontenc}
\usepackage[latin9]{inputenc}
\usepackage[toc,page]{appendix}
\usepackage{geometry}
\geometry{verbose,tmargin=1.5in,bmargin=1.5in,lmargin=1.5in,rmargin=1.5in}
\usepackage{bm}
\usepackage{amstext}
\usepackage{comment}
\usepackage{amsmath}
\usepackage{hyperref}
\hypersetup{
     colorlinks = true,
     citecolor = navy,
     linkcolor = navy,
     urlcolor=navy
}
\usepackage{adjustbox}
\usepackage[space]{grffile}
\usepackage{setspace}
\usepackage[normalem]{ulem}
\usepackage{nicematrix}
\usepackage{xcolor}
\usepackage{soul}
\usepackage{arydshln}
\usepackage{dsfont}
\usepackage{refstyle}
\usepackage[capitalize,nameinlink]{cleveref}
\usepackage{setspace}
\usepackage{bbm}
\usepackage[tableposition=t]{caption} % ensures caption of longtables is normalsized.
\usepackage{mathpazo} % add Palatino font
\usepackage{longtable}
\usepackage{amssymb}
\usepackage{amsfonts}
\usepackage{amsmath}
%\usepackage[margin=1in]{geometry} % to set margins
\usepackage{scalefnt}
\usepackage{wrapfig}
\usepackage{graphicx}
\usepackage{graphics}

%\usepackage{subfigure}
\usepackage{lscape}
\usepackage{pdflscape}
\usepackage{multirow}
\usepackage{array}
\usepackage{rotating}
\usepackage{bibunits}
\usepackage{fix-cm}
\usepackage{float}
\floatplacement{table}{H}
\usepackage{epsfig}
\usepackage{epstopdf}
\usepackage{refstyle}
\usepackage{tikz}
\usepackage{booktabs}
\usepackage{tabularx}
\usepackage{array}
\newcolumntype{x}[1]{>{\raggedleft\arraybackslash}p{#1}}%
\newcolumntype{d}[1]{>{\centering\arraybackslash}p{#1}}%
\newcolumntype{y}[1]{>{\raggedright\arraybackslash}p{#1}}%
\usepackage{verbatim}
%\epstopdfsetup{suffix=}
\setcounter{MaxMatrixCols}{10}
\usepackage[longnamesfirst,authoryear]{natbib}
\newenvironment{proof}[1][Proof]{\noindent \textbf{#1.}}{\rule{0.5em}{0.5em}}
\setlength{\topmargin}{-0.80in} \setlength{\textheight}{9.50in}
\setlength{\evensidemargin}{.0in} \setlength{\oddsidemargin}{.0in}
\setlength{\textwidth}{6.85in}
\setlength{\voffset}{0.4cm}
\setlength{\hoffset}{-0.5cm}
\usepackage{setspace}
\usepackage{afterpage}
\usepackage[flushleft]{threeparttable}
\usepackage{tabulary}
%\usepackage[titletoc,title]{appendix}
\usepackage[bottom]{footmisc}
\usepackage{color}
	\definecolor{maroon}{RGB}{128,0,0}
	\definecolor{navy}{RGB}{0,0,165}
	\definecolor{lb}{RGB}{150,165,190}
	\definecolor{bluenice}{rgb}{0.25,0.51,0.76}
	\definecolor{orangenice}{HTML}{EB811B}
	\definecolor{greennice}{HTML}{14B03D}
\usepackage{colortbl}
\usepackage{caption}
\usepackage{subcaption}
\usepackage{placeins}
\usepackage{nicefrac}
%\usepackage{chngcntr}
\newtheorem{assumption}{Assumption}
\usetikzlibrary{matrix,decorations.pathreplacing}

\newcommand{\sigmac}{\underline{\sigma}_{c,t}}
\newcommand{\sigmad}{\underline{\sigma}_{d,t}}
\newcommand{\dZ}{d\underline{Z}_t}
\newcommand{\sigmacd}{\underline{\sigma}_t^{\Delta}}
\newcommand{\D}{\mathcal{D}_t}
\newcommand{\sigmap}{\underline{\sigma}_{P,t}}
\newcommand{\sigmapc}{\underline{\sigma}_{P_c,t}}
\newcommand{\sigmapd}{\underline{\sigma}_{P_d,t}}
\newcommand{\mup}{\mu_t^{P}}
\newcommand{\mud}{\mu_t^{\mathcal{D}}}
\newcommand{\sigmaD}{\underline{\sigma}_{\mathcal{D},t}}
\newcommand{\muq}{\mu_{q,t}}
\newcommand{\sigmavt}{\underline{\sigma}_{\vartheta,t}}
\newcommand{\sigmaq}{\underline{\sigma}_{q,t}}
\newcommand{\sigmaK}{\underline{\sigma}_{K,t}}
\newcommand{\vt}{\vartheta_t}
\newcommand{\amb}{(\mathds{1}_a\sigma_a-\mathds{1}_b\sigma_b)}
\newcommand{\bma}{(\mathds{1}_b\sigma_b-\mathds{1}_a\sigma_a)}
\newcommand{\nuad}{\underline{\nu}_t^{ad}}
\newcommand{\nubd}{\underline{\nu}_t^{bd}}
\newcommand{\nuab}{\underline{\nu}_t^{a}}
\newcommand{\nubb}{\underline{\nu}_t^{b}}
\newcommand{\sigmaeta}{\underline{\sigma}_{\eta,t}}
\newcommand{\sigmas}{\underline{\sigma}_{S,t}}
\newcommand{\bsigmas}{\underline{\bar{\sigma}}_{S}}
\newcommand{\muvt}{\mu_t^\vartheta}
\newcommand{\PP}{\left(\underline{P}\right)}
\newcommand{\IO}{\mathbf{\Omega\left(\underline{P}\right)}}
\newcommand{\io}{\mathbf{\Omega}}
\newcommand{\fd}{\mathbf{\Xi\left(\underline{P}\right)}}
\newcommand{\fdd}{\mathbf{\Xi}}
\newcommand{\ishare}{\underline{\gamma}\left(\underline{P}\right)}
\newcommand{\ishares}{\underline{\gamma}}
\newcommand{\MM}{\mathbf{\mathcal{M}\PP}}
\newcommand{\MMM}{\mathbf{\mathcal{M}}}
\newcommand{\obj}{\mathbf{\Phi}}
\newcommand{\objt}{\mathbf{\tilde{\Phi}}}
\newcommand{\nuu}{(\underline{\nu})}
\newcommand{\nured}{\textcolor{red}{\nu}}
\newcommand{\nuredOne}{\textcolor{red}{\nu_1}}
\newcommand{\nuredTwo}{\textcolor{orangenice}{\nu_2}}
\newcommand{\scaleBlueOne}{\textcolor{bluenice}{\frac{1-\psi_{PHL,rice\rightarrow 11}\nu}{1-\psi_{PHL,1\rightarrow 11}}}}
\newcommand{\scaleBlueTwo}{\textcolor{bluenice}{\frac{1-\psi_{PHL,rice\rightarrow 1,rice}\nu}{1-\psi_{PHL,rice\rightarrow 1,rice}}}}
\newcommand{\scaleBlueOneOne}{\textcolor{bluenice}{\frac{1-\psi_{PHL,rice\rightarrow 11}\nu_1}{1-\psi_{PHL,1\rightarrow 11}}}}
\newcommand{\scaleBlueTwoOne}{\textcolor{bluenice}{\frac{1-\psi_{PHL,rice\rightarrow 1,rice}\nu_1}{1-\psi_{PHL,rice\rightarrow 1,rice}}}}

\newcommand{\scaleBlueOneTwo}{\textcolor{greennice}{1-\frac{1-\psi_{PHL,rice\rightarrow 11}}{\nu_2 }}}
\newcommand{\scaleBlueTwoTwo}{\textcolor{greennice}{1-\frac{1-\psi_{PHL,rice\rightarrow 1,rice}}{\nu_2}}}
\newcommand{\mb}[1]{\mathbold{\tilde{#1}}_t}
\newcommand{\mbb}[2]{\mathbold{\tilde{#1}}_{#2,t}}

	\definecolor{maroon}{RGB}{128,0,0}
	\definecolor{navy}{RGB}{0,0,165}
	\definecolor{lb}{RGB}{150,165,190}
	\definecolor{MyWhite}{RGB}{255, 255, 255}
	
\definecolor{red}{rgb}{0.9,0.1,0.1}
\definecolor{blue}{rgb}{0,0,0.7}
\definecolor{dgr}{rgb}{0.08,0.25,0.18}
\definecolor{navy}{RGB}{0,0,165}
\definecolor{org}{HTML}{EB811B}
\definecolor{lgr}{HTML}{14B03D}
\definecolor{shadecolor}{rgb}{0.8,0.8,0.8}

\definecolor{graphblue}{RGB}{64,130.56,194.5}
\definecolor{graphred}{RGB}{255,0,0}
\definecolor{graphyellow}{RGB}{237.8340,177.64,32}
\definecolor{graphpurple}{RGB}{126,47,142}

\newcommand{\hlight}[1]{\textcolor{org}{#1}}
\newcommand{\glight}[1]{\textcolor{lgr}{#1}}
\newcommand{\rlight}[1]{\textcolor{red}{#1}}
\newcommand{\blight}[1]{\textcolor{graphblue}{#1}}
\newcommand{\plight}[1]{\textcolor{graphpurple}{#1}}
\begin{document}
\title{Rethinking Supply and Demand Influence}
\author{N}
\maketitle
\vspace{-1.2cm}
\section{Purpose}
My apologies but I think there is a way easier way to do the influence calculations.

If this is correct, then we can save a bunch of time and make the implementation more transparent.

First I cover the logic and then the steps.

\section{The logic}

What we do is classify the \emph{value added} of a sector as having experienced a demand shock or a supply shock. Then what we want to do is to figure out the \emph{influence} of this on other sectors' prices.

There is a simple formula for this influence:

\[
	\underline{\log P}_t = (I-diag(\gamma)\mathbold{\Omega})^{-1}diag(1-\underline{\gamma})\log \underline{P}_{VA,t}
\]
where $\mathbold{\Omega}$ is the IO matrix whose row-sum is equal to 1 (taking a sum across column for a given row) and $\gamma$ is the share of sales that is spent on intermediates.

This formula states that the price in a given sector is a function of the value added prices in all sectors with weights given by $\mathcal{K} \equiv (I-\mathbold{\Omega})^{-1}diag(1-\underline{gamma})$. Each row of $\mathbold{\Theta}$ should sum to 1 over columns and gives the shares coming from each sector's value added price. Furthermore, each value should be positive.

We can calculate using this formula the share of each output price $\log P_i$ coming from the price of value added in sector $j$:

\[
	s_{j/i} = \mathbold{\Theta}_{i,j} \frac{P_{VA,j}}{P_i}
\]
where rows of $share_{j/i}$ are $i$ and columns are $j$. 

Then let $\underline{D}_t$ be an $N$ x 1 vector with entry $d_{j,t}=1$ if that sector's value added had a demand shock and $d_j=0$ if it had a supply shock. We can calculate the share of sector $i$'s \emph{output} price that is driven by supply or demand based on:

\[
	\underline{s}_{demand,t}  = \mathbold{s}\;\underline{D}_t
\]

That gives us the share of output prices in each sector that are driven by demand or supply.

Next, to do CPI, we simply multiply this by the CPI expenditure shares in time period $t$: $E_t$

\[
	share_{cpi,t} = \underline{E}_t' \;\underline{s}_{demand,t}
\]

\section{Coding this up}

Let's follow these simplified steps to get the decomposition done. The object we want is an $N$ x $T$ matrix $\rlight{\mathbold{D}}$ with entries $\{d_{i,t}\}$ between 0 and 1 that mean that sector $i$'s output prices in time $t$ are affected $d_{i,t}$ by demand shocks and $1-d_{i,t}$ by supply shocks. What we have is an $N$ x $T$ matrix $\hlight{\mathbold{\tilde{D}}}$ which is 1 whenever a sectors' value added receives a (net) demand shock and 0 when it receives a (net) supply shock. This is \underline{not} the same. $\hlight{\mathbold{\tilde{D}}}$ tells us if a demand shock \emph{began} in sector $i$ at time $t$ whereas  $\rlight{\mathbold{D}}$ tells us the \emph{net} effect of supply/demand shocks on sector $i$ and time $t$.

\begin{enumerate}
	\item Let's begin with  Let's collect all the (log of) value added prices calculated into a $N$ x $T$ matrix and call it $\glight{\mathbold{\log (P)}_{VA}}$. Finally collect together a $N$ x $T$ matrix of all log sectoral \emph{output} prices $\blight{\mathbold{\log(P)}}$.
	\item Then create an $N$ x $N$ matrix $\mathbold{\Theta}\equiv (I_N-diag(\underline{\gamma})\mathbold{\Omega})^{-1}diag(1-\underline{\gamma})$. Confirm the following:\footnote{Email me if anything doesn't look right :)}
	\begin{enumerate}
		\item All entries of $\mathbold{\Theta}$ are positive or 0.
		\item The sum of all rows (rowsum, sum across columns for a given row) equal 1
		\item That $\blight{\mathbold{\log(P)}} = \mathbold{\Theta} \glight{\mathbold{\log(P)}_{VA}}$.
	\end{enumerate}
	\item Once we have confirmed that $\mathbold{\Theta}$ is working as it should, let's now calculate value added prices conditional on there being a demand shock or a supply shock. Specifically, define:
	
	\begin{enumerate}
		\item $\mathbold{\log(P)_{\glight{VA}}^{(D=1)}} = \glight{\mathbold{\log (P)}_{VA}}\circ \hlight{\mathbold{\tilde{D}}}$ where $A\circ B$ denotes element multiplication. Note that this should make all entries of $\glight{\mathbold{\log(P)_{VA}}}$ 0 when $\hlight{\mathbold{\tilde{D}}}$ is 0. All other entries should be unchanged because $\hlight{\mathbold{\tilde{D}}}\in \{0,1\}$.
		\item $\mathbold{\log(P)_{\glight{VA}}^{(D=0)}} = \glight{\mathbold{\log (P)}_{VA}}\circ (1_{N,T}-\hlight{\mathbold{\tilde{D}}})$ where $1_{N,T}$ is an $N$ x $T$ matrix of ones.
	\end{enumerate}
	\item Calculate:
	\begin{enumerate}
		\item $\mathbold{\log(P)^{(D=1)}} = \mathbold{\Theta}\mathbold{\log(P)_{\glight{VA}}^{(D=1)}}$
		\item $\mathbold{\log(P)^{(D=0)}} = \mathbold{\Theta}\mathbold{\log(P)_{\glight{VA}}^{(D=0)}}$
	\end{enumerate}
	and confirm that $\mathbold{\log(P)^{(D=1)}} + \mathbold{\log(P)^{(D=0)}} = \blight{\mathbold{\log(P)}}$.
	\item Once confirmed we can calculate the final $\rlight{\mathbold{D}}$ matrix:
	\[
		\rlight{\mathbold{D}}\equiv \mathbold{\log(P)^{(D=1)}}\slash \blight{\mathbold{\log(P)}}
	\]
	where $\mathbold{A}\slash \mathbold{B}$ refers to element division (elements of $\mathbold{A}$ divided by elements of $\mathbold{B}$).
	\item Check that all elements of $\rlight{\mathbold{D}}$ are between 0 and 1.
	
\end{enumerate}


Contact me if anything is awry!


\section{An alternative way to calculate inflation}
Instead of calculating inflation as the percent change, why don't we calculate it as a log difference:

\[
	\pi_{i,t-1,t}\equiv \log(P_{i,t}) - \log(P_{i,t-1})
\]

For CPI inflation we have 

\[
	\pi_{i,t-1,t}^{CPI} \equiv \sum_i \omega_{i,t-1} ( \log(P_{i,t}) - \log(P_{i,t-1}))
\]

Once we know what share of sectors are supply and demand constrained we can decompose CPI into:

\[
	\pi_{i,t-1,t}^{CPI} = \underbrace{\sum_i \omega_{i,t-1}\mathds{1}_{D_{i,t}=1}\pi_{i,t-1,t}}_{\text{Demand Driven Inflation}} + \underbrace{\sum_i \omega_{i,t-1}\mathds{1}_{D_{i,t}=0}\pi_{i,t-1,t}}_{\text{Demand Driven Inflation}}  
\]

This is the easy bit but what happens if we want to calculate the 12 month moving average? With the log-method we can just sum the monthly changes:

\[
	\pi_{t-12,t}^{CPI} = \sum_{k=0}^{11} \pi_{t-k-1,t-k}^{CPI}
\]

Then we can do our decomposition super easily:

\begin{align*}
	\pi_{t-12,t}^{cpi} &= \sum_{k=0}^{11} \left(\sum_i \omega_{i,t-k-1}\mathds{1}_{D_{i,t-k}=1}\pi_{i,t-k-1,t-k}+\sum_i \omega_{i,t-k-1}\mathds{1}_{D_{i,t-k}=0}\pi_{i,t-k-1,t-k}\right)\\
	&=  \underbrace{\sum_i \sum_{k=0}^{11}\omega_{i,t-k-1}\mathds{1}_{D_{i,t-k}=1}\pi_{i,t-k-1,t-k}}_{\text{Demand driven inflation}}+\underbrace{\sum_i \sum_{k=0}^{11}\omega_{i,t-k-1}\mathds{1}_{D_{i,t-k}=0}\pi_{i,t-k-1,t-k}}_{\text{Supply driven inflation}}\\
\end{align*}

As with Shapiro's approach, we get a complete separation between demand and supply driven inflation but this approach has the advantage of the two types of inflation adding up to CPI exactly. 

The downside of this approach is that CPI inflation is not a percent change but a log-difference and these are similar but not identical. 

I don't have strong preferences either way.
\end{document}